% =========================================================
% Logical Strength and Quantifier Order
% =========================================================

\subsection*{Logical Strength and Quantifier Order}

\begin{remark*}[Section toolkit: Strength and order]
Quantifier order is one of the main new phenomena in predicate logic. In
propositional logic, changing the order of independent connectives often keeps
truth conditions recognizable. In predicate logic, moving \(\forall\) and
\(\exists\) can change whether witnesses may depend on earlier choices.
\end{remark*}

\begin{definitionbox}{Definition (Stronger Formula)}
\begin{definition}[Stronger Formula]
\label{def:strength}
A sentence \(A\) is \emph{stronger than} a sentence \(B\) exactly when
\[
A\models B
\quad\text{and}\quad
B\not\models A.
\]
\end{definition}
\end{definitionbox}

\begin{remark*}[Standard quantified statement]
\[
A \text{ is stronger than } B
\Longleftrightarrow
(\forall \mathcal M\;(\mathcal M\models A\Rightarrow\mathcal M\models B))
\land
(\exists \mathcal N\;(\mathcal N\models B\land\mathcal N\not\models A)).
\]
\end{remark*}

\begin{remark*}[Predicate reading]
\(\operatorname{StrongerFormula}(A,B)\) means that \(A\) entails \(B\), but
\(B\) does not entail \(A\).
\end{remark*}

\begin{remark*}[Negated quantified statement]
\[
A\not\models B
\quad\text{or}\quad
B\models A.
\]
\end{remark*}

\begin{remark*}[Interpretation]
The stronger sentence has fewer possible structures in which it is true. It
rules out more cases.
\end{remark*}

\begin{remark*}[Examples]
\(\forall x\,P(x)\) is stronger than \(\exists x\,P(x)\) over nonempty
domains.
\end{remark*}

\begin{dependencies}
\begin{itemize}
  \item \hyperref[def:log-consequence]{Logical Consequence}
\end{itemize}
\end{dependencies}

\begin{definitionbox}{Definition (Weaker Formula)}
\begin{definition}[Weaker Formula]
\label{def:weaker-formula}
A sentence \(B\) is \emph{weaker than} a sentence \(A\) exactly when \(A\) is
stronger than \(B\).
\end{definition}
\end{definitionbox}

\begin{remark*}[Standard quantified statement]
\[
B \text{ is weaker than } A
\Longleftrightarrow
A\models B
\text{ and }
B\not\models A.
\]
\end{remark*}

\begin{remark*}[Predicate reading]
\(\operatorname{WeakerFormula}(B,A)\) means that \(B\) is entailed by \(A\) but
does not entail \(A\).
\end{remark*}

\begin{remark*}[Interpretation]
The weaker sentence permits more structures. It follows from the stronger
sentence but carries less information.
\end{remark*}

\begin{dependencies}
\begin{itemize}
  \item \hyperref[def:strength]{Stronger Formula}
\end{itemize}
\end{dependencies}

\begin{definitionbox}{Definition (Quantifier Order)}
\begin{definition}[Quantifier Order]
\label{def:q-alt}
The \emph{quantifier order} of a prenex formula is the ordered list of
quantifiers in its prefix. A change of order is significant when it changes
which witnesses may depend on earlier variables.
\end{definition}
\end{definitionbox}

\begin{remark*}[Standard quantified statement]
\[
Q_1x_1\cdots Q_nx_n\,\psi
\text{ has quantifier order }
(Q_1,\ldots,Q_n).
\]
\end{remark*}

\begin{remark*}[Predicate reading]
\(\operatorname{QuantifierOrder}(\Phi,(Q_1,\ldots,Q_n))\) records the order of
the quantifiers in the prefix of \(\Phi\).
\end{remark*}

\begin{remark*}[Interpretation]
Quantifier order controls dependency. In \(\forall x\,\exists y\), the witness
\(y\) may depend on \(x\). In \(\exists y\,\forall x\), the same \(y\) must work
for all \(x\).
\end{remark*}

\begin{dependencies}
\begin{itemize}
  \item \hyperref[def:quantifier-prefix]{Quantifier Prefix}
  \item \hyperref[def:prenex-formula]{Prenex Formula}
\end{itemize}
\end{dependencies}

\begin{theorembox}{Theorem (Universal Implies Existential)}
\begin{theorem}[Universal Implies Existential]
\label{thm:universal-implies-existential}
Over nonempty domains,
\[
\forall x\,\varphi(x)\models \exists x\,\varphi(x).
\]
\hyperref[prf:universal-implies-existential]{Go to proof.}
\end{theorem}
\end{theorembox}

\begin{remark*}[Standard quantified statement]
\[
\forall \mathcal M\;
\bigl(\mathcal M\models\forall x\,\varphi(x)
\Rightarrow
\mathcal M\models\exists x\,\varphi(x)\bigr).
\]
\end{remark*}

\begin{remark*}[Predicate reading]
\(\operatorname{LogicalConsequence}(\{\forall x\,\varphi(x)\},\exists x\,\varphi(x))\).
\end{remark*}

\begin{remark*}[Interpretation]
If every domain element satisfies \(\varphi\), then at least one does because
first-order structures have nonempty domains.
\end{remark*}

\begin{dependencies}
\begin{itemize}
  \item \hyperref[def:structure]{Structure}
  \item \hyperref[def:log-consequence]{Logical Consequence}
\end{itemize}
\end{dependencies}

\begin{theorembox}{Theorem (Existential Does Not Imply Universal)}
\begin{theorem}[Existential Does Not Imply Universal]
\label{thm:existential-not-universal}
In general,
\[
\exists x\,\varphi(x)\not\models \forall x\,\varphi(x).
\]
\hyperref[prf:existential-not-universal]{Go to proof.}
\end{theorem}
\end{theorembox}

\begin{remark*}[Standard quantified statement]
\[
\exists \mathcal M\;
\bigl(\mathcal M\models\exists x\,\varphi(x)
\land
\mathcal M\not\models\forall x\,\varphi(x)\bigr).
\]
\end{remark*}

\begin{remark*}[Predicate reading]
\(\neg\operatorname{LogicalConsequence}(\{\exists x\,\varphi(x)\},\forall x\,\varphi(x))\).
\end{remark*}

\begin{remark*}[Interpretation]
One witness does not make a universal claim true.
\end{remark*}

\begin{remark*}[Exposition]
In the real numbers, \(\exists x\,(x=0)\) is true, but
\(\forall x\,(x=0)\) is false.
\end{remark*}

\begin{dependencies}
\begin{itemize}
  \item \hyperref[thm:universal-implies-existential]{Universal Implies Existential}
\end{itemize}
\end{dependencies}

\begin{theorembox}{Theorem (Uniform Witness Implies Dependent Witness)}
\begin{theorem}[Uniform Witness Implies Dependent Witness]
\label{thm:uniform-witness-implies-dependent-witness}
For every formula \(\varphi(x,y)\),
\[
\exists y\,\forall x\,\varphi(x,y)
\models
\forall x\,\exists y\,\varphi(x,y).
\]
\hyperref[prf:uniform-witness-implies-dependent-witness]{Go to proof.}
\end{theorem}
\end{theorembox}

\begin{remark*}[Standard quantified statement]
\[
\forall \mathcal M\;
\bigl(
\mathcal M\models\exists y\,\forall x\,\varphi(x,y)
\Rightarrow
\mathcal M\models\forall x\,\exists y\,\varphi(x,y)
\bigr).
\]
\end{remark*}

\begin{remark*}[Predicate reading]
\(\operatorname{LogicalConsequence}(\{\exists y\,\forall x\,\varphi(x,y)\},
\forall x\,\exists y\,\varphi(x,y))\).
\end{remark*}

\begin{remark*}[Interpretation]
A single witness \(y\) that works for every \(x\) can be reused as the
witness required for each separate \(x\).
\end{remark*}

\begin{dependencies}
\begin{itemize}
  \item \hyperref[def:q-alt]{Quantifier Order}
  \item \hyperref[def:log-consequence]{Logical Consequence}
\end{itemize}
\end{dependencies}

\begin{theorembox}{Theorem (Dependent Witness Need Not Be Uniform)}
\begin{theorem}[Dependent Witness Need Not Be Uniform]
\label{thm:dependent-witness-not-uniform}
In general,
\[
\forall x\,\exists y\,\varphi(x,y)
\not\models
\exists y\,\forall x\,\varphi(x,y).
\]
\hyperref[prf:dependent-witness-not-uniform]{Go to proof.}
\end{theorem}
\end{theorembox}

\begin{remark*}[Standard quantified statement]
\[
\exists \mathcal M\;
\bigl(
\mathcal M\models\forall x\,\exists y\,\varphi(x,y)
\land
\mathcal M\not\models\exists y\,\forall x\,\varphi(x,y)
\bigr).
\]
\end{remark*}

\begin{remark*}[Predicate reading]
\(\neg\operatorname{LogicalConsequence}(\{\forall x\,\exists y\,\varphi(x,y)\},
\exists y\,\forall x\,\varphi(x,y))\).
\end{remark*}

\begin{remark*}[Interpretation]
The dependent-witness statement allows \(y\) to vary with \(x\). The uniform
statement requires one \(y\) to handle all \(x\) at once.
\end{remark*}

\begin{remark*}[Exposition]
Over \(\mathbb R\), let \(\varphi(x,y)\) be \(y>x\). Then
\(\forall x\,\exists y\,(y>x)\) is true, since \(y=x+1\) works for each \(x\).
But \(\exists y\,\forall x\,(y>x)\) is false, since no real number is larger
than every real number.
\end{remark*}

\begin{remark*}[Exposition]
Quantifier order matters because it records dependency. The phrase "for every
\(x\), there exists \(y\)" permits a different \(y\) after \(x\) is known. The
phrase "there exists \(y\) such that for every \(x\)" chooses \(y\) before
seeing \(x\). This is a genuine expressive jump beyond propositional logic.
\end{remark*}

\begin{dependencies}
\begin{itemize}
  \item \hyperref[thm:uniform-witness-implies-dependent-witness]{Uniform Witness Implies Dependent Witness}
\end{itemize}
\end{dependencies}
